\section{Ethics, fairness and privacy}
\label{sec:ethics}

\paragraph{Who bears the errors.} A false positive in this system declines or holds a legitimate
payment. Which customers bear those errors is not established here: the specification requires a
fairness report
(in its evaluation protocol): false-positive rate and blocked legitimate amount by country, channel,
KYC tier and urban or rural location. That report is \notyet{M4 follow-up}. The per-country and
per-channel AUCs in Sections~\ref{sec:evaluation} and~\ref{sec:generalisation} are not a fairness
analysis: they measure ranking, not who is wrongly blocked. On a synthetic benchmark whose
population structure is partly sourced and whose fraud is assumed, even the fairness report will
describe the generator's disparities, not a real population's.

\paragraph{Customer verification.} The specification's self-service unblocking by SMS link would
be defeated by the fraud it targets: in a SIM-swap takeover the attacker holds the victim's SIM and
receives the link. Self-service is therefore specified to be disabled when the last SIM swap is less than
\nSelfServiceSimDays{} days old or its recency is unknown, when the device changed in the last
\nSelfServiceDeviceHours{} hours, when the ensemble score is at least \nSelfServiceScore{}, or when fraud-type
signals indicate account takeover, and the customer is directed to the institution's official number (resolution D-25). This is specified behaviour of
a later milestone, not an evaluated result.

\paragraph{Privacy.} The benchmark contains no real personal data. The system is specified to keep
personal data out of every analytic path: account and counterparty identifiers must be opaque tokens
and raw phone and account numbers are rejected at ingestion (the specification's ingestion
validation), personal data lives only in a separate encrypted vault (resolution D-20), and logs pass
through redacting processors (the specification's observability requirements). These are specified
behaviours; this paper does not evaluate them.
Explanations shown to analysts and customers are specified to come from deterministic
templates, never from a language model; those templates have not been written, and the analyst
console that would show explanations is not part of this paper's release (Section~\ref{sec:method}).

\paragraph{Regulation and residency.} Statements about reporting obligations and data protection
in Rwanda (for example the data protection law and the financial intelligence reporting framework)
are unverified and must be confirmed by a compliance professional (resolution D-22); the API
contract defines a marker for report fields that are unconfirmed, and no report template exists. Production data must reside in a facility that is
legally permitted for it; the infrastructure code cannot satisfy that requirement, and
demonstration environments run on synthetic data only, labelled as such (resolution D-21).

\paragraph{What this paper does not claim.} That FraudShield is deployed at, endorsed by, or
validated with any real institution or regulator, or that any figure here predicts performance on
real transactions.
